\begin{abstract}
Kaladont is a word chain game, popular among children in Serbo-Croatian speaking regions, where each new word must start with the last two letters of the previous one.
The game ends when a player cannot find a valid word to continue the chain. In this paper we investigate the feasibility of playing the game of Kaladont in different languages. We conjecture that languages that possess the property of phonemic orthography, where the written text directly mirrors the spoken language — one letter represents exactly one sound (phoneme), and every letter is pronounced, are better suited for playing the game. To find out, we randomly simulate a large number of games in many languages from different language families, using  word lists from online data sources.
Our results show average chain lengths range from 2 words to 48 words across different languages, and that, even though Serbo-Croatian comes out on top, yes, there are languages that come very close and for which playing the game of Kaladont is very feasible.
\end{abstract}
